\chapter{Frontend Auth State (AuthContext)}

\section{Why a Context Instead of Prop Drilling}

Whether a user is logged in, and who they are, is needed by components
scattered across the whole tree --- the navbar shows their name, protected
routes need to check they're logged in at all, and profile pages need their
data. Passing that down as props through every intermediate component would
mean threading auth state through components that don't care about it at
all. React's Context API lets any component subscribe directly to auth
state without that chain.

\begin{diagrambox}[AuthProvider Tree]
\begin{verbatim}
              <AuthProvider>
                     |
      -----------------------------------
      |            |             |      |
  <Navbar>   <Dashboard>    <Profile>  <ProtectedRoute>
   (shows      (shows         (shows     (checks
   user.name)  user's tasks)  user data) currentUser)
\end{verbatim}
\end{diagrambox}

\section{AuthContext Implementation}

\begin{lstlisting}[language=JavaScript]
// frontend/src/context/AuthContext.jsx
import { createContext, useContext, useState, useEffect } from "react";
import api from "../services/api";

const AuthContext = createContext(null);

export function AuthProvider({ children }) {
  const [currentUser, setCurrentUser] = useState(null);
  const [loading, setLoading] = useState(true);

  // Restore session on app load by asking the backend who the
  // current cookie/token belongs to, if anyone.
  useEffect(() => {
    const checkAuth = async () => {
      try {
        const { data } = await api.get("/auth/me");
        setCurrentUser(data.user);
      } catch {
        setCurrentUser(null);
      } finally {
        setLoading(false);
      }
    };
    checkAuth();
  }, []);

  const login = async (credentials) => {
    const { data } = await api.post("/auth/login", credentials);
    setCurrentUser(data.user);
    return data.user;
  };

  const logout = async () => {
    await api.post("/auth/logout");
    setCurrentUser(null);
  };

  const value = { currentUser, loading, login, logout };

  return <AuthContext.Provider value={value}>{children}</AuthContext.Provider>;
}

export function useAuth() {
  const context = useContext(AuthContext);
  if (!context) {
    throw new Error("useAuth must be used within an AuthProvider");
  }
  return context;
}
\end{lstlisting}

\begin{notebox}[NOTE]
\texttt{loading} exists so consumers (especially
\texttt{ProtectedRoute}) can distinguish ``we don't know yet whether the
user is logged in'' from ``we checked, and they are not.'' Without it, a
protected route would briefly redirect to \texttt{/login} on every page
refresh before \texttt{/auth/me} has had a chance to respond.
\end{notebox}

\section{Wiring the Provider in main.jsx}

\begin{lstlisting}[language=JavaScript]
// frontend/src/main.jsx
import React from "react";
import ReactDOM from "react-dom/client";
import { BrowserRouter } from "react-router-dom";
import App from "./App";
import { AuthProvider } from "./context/AuthContext";

ReactDOM.createRoot(document.getElementById("root")).render(
  <React.StrictMode>
    <BrowserRouter>
      <AuthProvider>
        <App />
      </AuthProvider>
    </BrowserRouter>
  </React.StrictMode>
);
\end{lstlisting}

\section{Consuming the Context: Navbar}

\begin{lstlisting}[language=JavaScript]
// frontend/src/components/Navbar.jsx
import { useNavigate } from "react-router-dom";
import { useAuth } from "../context/AuthContext";

function Navbar() {
  const { currentUser, logout } = useAuth();
  const navigate = useNavigate();

  const handleLogout = async () => {
    await logout();
    navigate("/login");
  };

  return (
    <nav>
      <span className="brand">Task Manager</span>
      {currentUser ? (
        <div className="navbar-user">
          <span>Hi, {currentUser.name}</span>
          <button onClick={handleLogout}>Logout</button>
        </div>
      ) : (
        <a href="/login">Login</a>
      )}
    </nav>
  );
}

export default Navbar;
\end{lstlisting}

\begin{tipbox}[TIP]
Every component that needs auth state --- \texttt{Navbar},
\texttt{ProtectedRoute}, \texttt{Profile} --- calls the same
\texttt{useAuth()} hook instead of importing \texttt{AuthContext} directly.
It's a small wrapper, but it gives a single place to add the
``used outside provider'' guard and keeps consumers decoupled from the
context's internal shape.
\end{tipbox}
